% =============================================================
% CONCLUSION
% =============================================================
\section{Conclusion}
\label{sec:conclusion}

This paper has constructed the first annual, internationally comparable series of Latvian GDP per capita for the Soviet period, 1937--95, from physical output, transport and employment data anchored to Bolotin's 1990 benchmark, and joined it to the published interwar series, extended to 1919 and 1940, and to adjusted modern data, to give a continuous series for 1913 and 1919--2024 in 1990 and 2011 international dollars. Built from different sources, with different weights and from benchmarks fifty-five years apart, the interwar and Soviet-period series agree within 3.5 per cent in every year of their 1937--40 overlap, and within 2 per cent in three of the four, and the Soviet-period series meets a modern series derived from the 2011 ICP round within 3 per cent at the 1990--1 splice. The series supersedes the Maddison Project's figures, which carry a benchmark inflated by the first year of transition backwards with net material product growth and so place Latvian income in 1990 a quarter too high while understating it in the 1970s, and it is ready for the Maddison Project Database. With the Lithuanian and Estonian series it completes the national accounts of the Baltic states and provides, for the first time, modern GDP estimates for three republics of the Soviet Union; the procedure can be applied to the others.

The Latvian series tells a story of inherited strength spent. Latvia lost more than a third of its population and three-fifths of its output in the First World War and spent the 1920s regaining its 1913 income; it grew by \gr{1929}{1938} per cent a year in the 1930s behind import-substituting industrialisation, lost a quarter of its population and 28 per cent of its income per head in the Second World War, recovered by \lvREGAINthirtyeight{}, and grew by \gr{1950}{1968} per cent a year between 1950 and 1968 as a machine-building and electrotechnical workshop of the Soviet Union, built with immigrant labour on the skills and infrastructure of the interwar republic \parencites[115--16]{vaskela_2013_lietuvos}[184]{years_of_dependence}[42--3]{puur_2019_turbulent}. From 1968 it stagnated earlier and more completely than its neighbours, at \gr{1968}{1989} per cent a year to 1989, as investment was redirected, construction contracted, harvests fell and the all-Union industries aged; it ceded the Baltic lead to Estonia in 1989, suffered the deepest transitional collapse of the three Baltic states, regained its 1989 income only in \lvREGAINeightynine{}, and was overtaken by Lithuania in \ltOVERTAKElv{}.

The comparison of the three Baltic economies across three regimes yields four conclusions. First, under interwar capitalism the region recovered rather than converged: the dynamic 1930s only restored the relative position of 1913, and the parity with Finland that Baltic memory assumes had already been lost by 1938 \parencites[228]{vilpisauskas_2014_lithuania}[154]{norkus_2023_baltic}. Second, Soviet socialism delivered two decades of genuine catch-up, faster than the Western core though slower than its periphery, through coercive mobilisation of labour and capital on a base of inherited social capability; and then two decades in which the region grew more slowly than any other in Europe, as the mechanisms of extensive growth were exhausted and the system proved unable to replace them. The gains of 1950--68 were cancelled by 1989, and the Baltic republics, among the best-placed peripheral economies of Europe in 1938, were poorer than Greece in 1989. Third, restored independence has brought the only sustained convergence of the century, at a demographic cost without parallel in post-socialist Europe, and it has been conditional on reform and uneven in outcome: Lithuania has exceeded the convergence expectation of the European pool, Estonia has exceeded it by less, and Latvia, which in 1938 was best placed of the three, has merely met it. Fourth, in the convergence regressions of each period no Baltic economy exceeded the expectation set by the European pool before 1989, and all three fell below it under socialism; against the Nordic countries none has regained its relative position of 1913.

The puzzle with which the paper began has, then, a partial answer. The Baltic states did not fail to converge for want of social capability; they possessed its human-capital component throughout, and lost half a century in which it could not be used, first in the autarkic and unstable environment of the interwar period, which limited what egalitarian and literate agrarian societies could achieve, and then in a system that could exploit their capability for a generation but not renew it. What they have made of their third opportunity differs, and the divergence since 2008 between Lithuania and its northern neighbours, which the Latvian series brings into focus, is the question the comparison leaves open. The series presented here pose that question, and the Soviet-era turning points of 1955 and 1968, with a precision that was not available before. Answering them will require what the series cannot provide: republican investment data corrected for Soviet pricing, of the kind Vonyó and Klein have built for East-Central Europe \parencite{vonyo_whydidsocialistfail}; a Soviet-period benchmark with republican rather than Union-wide prices; and consumption-side estimates against which the output-side record of Soviet living standards can be tested.
